\documentclass[10pt,twocolumn]{article}
\usepackage[T1]{fontenc}
\usepackage[utf8]{inputenc}
\usepackage{lmodern}
\usepackage[margin=1.8cm,top=2.4cm,bottom=2cm,columnsep=0.6cm]{geometry}
\usepackage{fancyhdr}
\usepackage{titlesec}
\usepackage{booktabs}
\usepackage{amsmath,amssymb,amsfonts}
\usepackage{microtype}
\usepackage{xcolor}
\usepackage{mdframed}
\usepackage{parskip}
\usepackage{enumitem}
\usepackage{hyperref}

\definecolor{rulered}{RGB}{180,30,30}
\definecolor{boxgray}{RGB}{245,245,245}
\definecolor{keywordgray}{RGB}{100,100,100}

\pagestyle{fancy}
\fancyhf{}
\fancyhead[L]{\textsc{\small Claude Mag}}
\fancyhead[C]{\small\textit{Machine Learning Research Digest}}
\fancyhead[R]{\small 2026-08-24}
\fancyfoot[C]{\small\thepage}
\renewcommand{\headrulewidth}{0.8pt}

\titleformat{\section}{\normalsize\bfseries\color{rulered}}{}{0pt}{}%
  [\vspace{-4pt}\color{rulered}\rule{\columnwidth}{0.4pt}]
\titlespacing{\section}{0pt}{8pt}{4pt}

\hypersetup{colorlinks=true,urlcolor=rulered,linkcolor=black}

\begin{document}
\setlength{\parindent}{0pt}
\setlength{\parskip}{4pt}

{\small\color{keywordgray}\textbf{Keywords:}
  VLM evaluation \textbullet{}
  behavioral reliability \textbullet{}
  scientific figure understanding \textbullet{}
  confabulation}\par\vspace{4pt}

{\LARGE\bfseries\raggedright
  How Do VLMs Behave When Blind or Misled?
  Behavioral Evaluation of VLMs on Scientific
  Figures}\par\vspace{4pt}

{\large\itshape\raggedright
  SciFigBench Reveals That VLMs Confabulate Data Values When Figures
  Are Blurred and Defer to Incorrect Captions Over Visual
  Evidence}\par\vspace{6pt}

{\small\textbf{By} Paul Osemudiame Oamen, Owusu-Banahene Osei,
  Ananya Mukherjee, Christian Greisinger, Steffen Eger,
  Pius Onobhayedo, Wei Zhao}\par\vspace{2pt}

{\small\color{keywordgray}arXiv:2608.13267 \textbullet{} Preprint, August 2026}
\par\vspace{2pt}\noindent{\color{rulered}\rule{\columnwidth}{0.6pt}}\vspace{6pt}

AI assistants for scientific reading increasingly interpret research
figures, but standard benchmarks test only accuracy on clean images.
SciFigBench stress-tests VLMs on 250 scientific figures across 34,000+
evaluation setups probing perception, reasoning, and behavioral
reliability under uncertainty. The findings are alarming: VLMs
confabulate specific data values in 62--84\% of cases when critical
figure elements are blurred, and defer to incorrect captions over
accurate visual evidence 41--67\% of the time. Models ranked highest
on standard accuracy correlate poorly with behavioral reliability
($\rho = 0.31$), indicating that safe scientific deployment cannot be
inferred from existing evaluations.

\begin{mdframed}[backgroundcolor=boxgray,linecolor=rulered,linewidth=1pt,
    innerleftmargin=6pt,innerrightmargin=6pt,
    innertopmargin=6pt,innerbottommargin=6pt]
\textbf{\small AT A GLANCE}\par\vspace{4pt}
\begin{description}[leftmargin=0pt,labelsep=4pt,itemsep=2pt,parsep=0pt,
    font=\small\bfseries]
  \item[Problem:] Standard VLM benchmarks measure accuracy on clean
    figures but not behavioral reliability when visual evidence is
    missing or misleading --- a critical gap for scientific AI
    assistants.
  \item[Why it matters:] Scientific AI tools are being deployed for
    literature review and data extraction; a model that fabricates
    values for blurred axes or trusts wrong captions poses real risks.
  \item[Method:] 250 expert-annotated figures extended via image
    transformations, resistance probes, caption-bias probes, and
    selective-blur targets, producing 34,000+ evaluation setups.
  \item[Result:] VLMs confabulate on 62--84\% of blur conditions;
    caption-bias rate 41--67\%; standard accuracy rank correlates
    poorly with behavioral reliability ($\rho = 0.31$).
  \item[Evidence:] Resistance probe failures (71\% compliance
    with false premises) confirm that models do not systematically
    detect their own visual limitations.
\end{description}
\end{mdframed}

\section{The Big Picture}

AI assistants for scientific literature are arriving rapidly: systems
that can summarize papers, extract data from figures, answer questions
about experimental results, and assist peer reviewers. These systems
rely on VLMs that have demonstrated strong performance on standard
visual question answering benchmarks.

Standard accuracy evaluations present VLMs with clean, fully-readable
figures and ask questions with unambiguous ground-truth answers. This
measures one important capability --- accurate perception and reasoning
--- but ignores a second capability that is equally important for
scientific use: \emph{knowing what you cannot see}.

When a research figure has a blurred axis label, a partially occluded
data series, or a misleading caption inserted during an adversarial
or unintentional editing step, a reliable scientific AI assistant
should express appropriate uncertainty. Instead, SciFigBench finds
that current VLMs confidently fabricate specific data values, comply
with false premises, and defer to textual context over visual evidence.

\section{Why Existing Approaches Fall Short}

\textbf{Standard VQA and figure understanding benchmarks} evaluate
on clean images with unambiguous ground truth. They do not assess:
\begin{itemize}[itemsep=2pt,parsep=0pt]
  \item Whether models appropriately express uncertainty when evidence
    is missing (rather than guessing)
  \item Whether models are biased by captions that contradict visual
    content
  \item Whether models detect fine-grained degradation of specific
    visual elements (e.g., one axis blurred while another is intact)
\end{itemize}

\textbf{Hallucination benchmarks} typically test factual grounding
in natural images (``is there a cat in the image?'') rather than
quantitative data extraction from structured scientific visualizations.

SciFigBench fills the gap by systematically varying the \emph{state}
of the visual evidence rather than just the questions asked about it.

\section{Core Mechanism}

\textbf{Base figures:} 250 scientific figures collected from STEM
research papers across biology, physics, chemistry, computer science,
and medicine. Each figure receives expert annotations identifying:
key data values, axis labels, trend directions, visual element
identities, and figure captions.

\textbf{Evaluation setups (>34,000 total):}
\begin{enumerate}[itemsep=2pt,parsep=0pt,label=(\arabic*)]
  \item \textbf{Image transformations:} Gaussian blur, label removal,
    color-region masking applied to specific elements while leaving
    others intact.
  \item \textbf{Reasoning questions:} Quantitative (``What is the
    value at $x=3$?'') and qualitative (``Which series shows the
    steeper slope?'').
  \item \textbf{Resistance probes:} Model is told (falsely) that
    a specific value can be read from a blurred region --- tests
    whether it complies or refuses.
  \item \textbf{Caption-bias probes:} Figure paired with a caption
    that states a trend opposite to what the graph shows --- tests
    whether model trusts visual evidence or caption.
  \item \textbf{Selective-blur targets:} Single data elements blurred
    while the rest of the figure is intact --- tests fine-grained
    awareness of local information availability.
\end{enumerate}

\section{Technical Formulation}

The \textbf{behavioral reliability score (BRS)} for model $M$ on
perturbation condition set $\mathcal{P}$:
\[
  \text{BRS}(M) = \frac{1}{|\mathcal{P}|}\sum_{c \in \mathcal{P}}
    \frac{1}{|Q_c|}\sum_{q \in Q_c}
    \mathbf{1}[M(f_c, q) \text{ is behaviorally appropriate}]
\]
where $f_c$ is figure $f$ under perturbation $c$, $Q_c$ is the
question set for condition $c$, and ``behaviorally appropriate''
means the model expresses uncertainty or refuses when evidence
is unavailable, and trusts visual evidence over conflicting captions.

The \textbf{caption-bias score (CBS)}:
\[
  \text{CBS}(M) = \mathbb{E}_{(f,q)}\!\left[\mathbf{1}\!\left[
    M(f, q, \text{cap}_\text{wrong}) \neq M(f, q, \text{no cap})
  \right]\right]
\]
where $\text{cap}_\text{wrong}$ is a caption that contradicts visual
truth.

The \textbf{refusal rate (RR)} on resistance probes:
\[
  \text{RR}(M) = \mathbb{E}_{(f,q,\text{probe})}\!\left[
    \mathbf{1}[M \text{ declines to provide value}]
  \right]
\]

\section{Learning or Inference Procedure}

SciFigBench is an evaluation benchmark with no training component.
Models are evaluated zero-shot with a standardized system prompt
that instructs them to answer to the best of their ability. No
special uncertainty-elicitation instructions are provided in the
main evaluation; ablations with explicit uncertainty instructions
show modest improvement but do not eliminate confabulation.

\section{What the Guarantee Says}

No formal guarantee is provided. Ground truth annotations were
produced by domain experts (PhD students and faculty in the
respective fields) with inter-annotator agreement verified via
Cohen's kappa ($\kappa > 0.82$ across all annotation types).
Behavioral appropriateness judgments are rule-based (not LLM-judge)
wherever possible: refusal is detected by keyword matching over
standardized response categories.

\section{Experimental Findings}

\begin{itemize}[itemsep=2pt,parsep=0pt]
  \item \textbf{Confabulation under blur:} VLMs provide specific
    numerical answers 62--84\% of the time when the asked-about
    visual element is blurred, nearly unchanged from clean conditions.
  \item \textbf{Caption bias:} Models align with the contradictory
    caption over correct visual evidence 41--67\% of the time.
    Larger models show more bias, not less.
  \item \textbf{Resistance probe compliance:} Models comply with false
    premises (generate values from blurred regions) 71\% of the time.
    Refusal rate is below 15\% for all evaluated models.
  \item \textbf{Rank-reliability correlation:} Spearman $\rho = 0.31$
    between standard accuracy rank and BRS rank, indicating that
    accuracy evaluations do not predict behavioral reliability.
  \item \textbf{Domain variation:} Biology and chemistry figures show
    higher caption-bias rates than computer science figures, possibly
    due to differences in training data figure-caption co-occurrence.
\end{itemize}

\section{Ablations and Interpretation}

\textbf{Uncertainty elicitation:} Prompting models to ``express
uncertainty when you cannot read a value'' reduces confabulation
from 74\% to 52\% --- improvement, but not elimination. Models still
confabulate more than half the time despite explicit instruction.

\textbf{Blur severity:} Confabulation rate is surprisingly stable
across blur intensities from mild ($\sigma=2$) to severe ($\sigma=20$).
This suggests models are generating values from semantic expectations
rather than partially visible pixel information.

\textbf{Caption-bias by model size:} The largest evaluated models
(7B--70B parameters) exhibit more caption bias than smaller models
(1B--3B), consistent with larger models placing more weight on
linguistic context relative to visual evidence during generation.

\textbf{Selective vs.\ global blur:} When only a single axis label
is blurred while the rest of the figure is intact, confabulation
rate (78\%) slightly exceeds the global-blur rate (74\%), suggesting
models do not notice fine-grained local information loss.

\section*{Reference}

Paul Osemudiame Oamen, Owusu-Banahene Osei, Ananya Mukherjee,
Christian Greisinger, Steffen Eger, Pius Onobhayedo, Wei Zhao.
\textbf{How Do VLMs Behave When Blind or Misled? Behavioral
Evaluation of VLMs on Scientific Figures.}
arXiv:2608.13267, August 2026.
\url{https://arxiv.org/abs/2608.13267}

\end{document}
